\documentclass[11pt]{article}

\usepackage[margin=0.85in]{geometry}
\usepackage{amsmath,amssymb,amsthm}
\usepackage{array}
\usepackage{booktabs}
\usepackage{xcolor}
\usepackage{hyperref}
\usepackage{url}
\usepackage{enumitem}
\usepackage{titlesec}

\definecolor{nicmeBlue}{HTML}{1F4E79}
\definecolor{nicmeTeal}{HTML}{1B7F79}
\definecolor{nicmeGold}{HTML}{B7791F}
\definecolor{nicmeLight}{HTML}{F4F7FA}
\definecolor{nicmeLine}{HTML}{D7DEE8}
\definecolor{nicmeText}{HTML}{1F2933}
\definecolor{nicmeWarn}{HTML}{7A4F01}

\hypersetup{
  colorlinks=true,
  linkcolor=nicmeBlue,
  urlcolor=nicmeTeal,
  citecolor=nicmeBlue,
  pdftitle={Why NICME Can Outperform CSADA on PMI-10},
  pdfauthor={NICME Repository Notes},
  pdfsubject={Cost-sensitive logit adjustment, expected-cost regularization, and CSADA},
  pdfkeywords={NICME, CSADA, cost-sensitive learning, adversarial data augmentation, logit adjustment, PMI}
}

\titleformat{\section}
  {\large\bfseries\color{nicmeBlue}}
  {\thesection}{0.7em}{}
\titleformat{\subsection}
  {\normalsize\bfseries\color{nicmeTeal}}
  {\thesubsection}{0.7em}{}
\setlist[itemize]{leftmargin=1.35em, itemsep=0.18em}
\setlist[enumerate]{leftmargin=1.45em, itemsep=0.18em}
\emergencystretch=3em

\newcommand{\code}[1]{\texttt{#1}}
\newcommand{\pathref}[1]{\texttt{\small #1}}
\newcommand{\softmax}{\operatorname{softmax}}
\newcommand{\argmax}{\operatorname*{arg\,max}}
\newcommand{\argmin}{\operatorname*{arg\,min}}
\newcommand{\E}{\mathbb{E}}
\newcommand{\R}{\mathbb{R}}
\newcommand{\I}{\mathbf{1}}
\newcommand{\keybox}[1]{%
  \vspace{0.35em}
  \noindent\fcolorbox{nicmeLine}{nicmeLight}{%
    \begin{minipage}{0.965\linewidth}
      #1
    \end{minipage}
  }
  \vspace{0.35em}
}
\newcommand{\warnbox}[1]{%
  \vspace{0.35em}
  \noindent\fcolorbox{nicmeGold}{white}{%
    \begin{minipage}{0.965\linewidth}
      \color{nicmeWarn}#1
    \end{minipage}
  }
  \vspace{0.35em}
}

\theoremstyle{definition}
\newtheorem{definition}{Definition}
\theoremstyle{plain}
\newtheorem{proposition}{Proposition}
\newtheorem{lemma}{Lemma}
\theoremstyle{remark}
\newtheorem{remark}{Remark}

\begin{document}

\begin{center}
  {\huge\bfseries\color{nicmeBlue} Why NICME Can Outperform CSADA}\\[0.25em]
  {\Large A Cost-Sensitive Margin And Expected-Cost View}\\[0.6em]
  {\normalsize Repository note for \pathref{/mnt/storage/github/NICME}}\\
  {\normalsize Current as of May 2026}
\end{center}

\vspace{0.55em}

\keybox{
\textbf{Thesis.}
For the completed PMI-10 balanced no-calibration experiment, NICME can plausibly outperform CSADA on the clean-test, recall-first, cost-sensitive tradeoff because its loss directly modifies the clean logits in two aligned ways: it imposes pairwise raw-logit margins proportional to the expert cost matrix, and it directly penalizes clean softmax probability mass on costly wrong classes.
CSADA is powerful, and was designed exactly to make overparameterized deep networks cost-sensitive, but its signal is mediated through local targeted perturbations and adversarial robust optimization.
That mediation can be less aligned with clean expected cost and can trade away clean balanced accuracy.
}

\warnbox{
\textbf{Scope of the claim.}
This is not a universal theorem that NICME dominates CSADA.
It is a theoretical and empirical explanation for the observed repository result:
NICME run 95 ties the previous CSADA result on target-min and target-macro recall, but has much lower normalized ATC and much higher balanced accuracy.
The strict validation-selected HPO row remains run 89, while run 95 is the best observed test row.
}

\section{Formal Problem Setup}

Let $(X,Y)\sim P$ with labels $Y\in[K]=\{1,\ldots,K\}$.
Let $\eta_y(x)=P(Y=y\mid X=x)$ be the conditional class probability and let
$C\in\R_{\ge 0}^{K\times K}$ be a fixed cost matrix with convention
$C_{y,k}=C[\text{true}=y,\text{pred}=k]$.
In this repository, correct predictions have cost $C_{y,y}=0$; ordinary noncritical off-diagonal PMI-10 errors have cost $1$; and four directed critical pairs have costs $8$ or $10$.

\begin{definition}[Cost risk and ATC]
For a classifier $h:\mathcal{X}\to[K]$, the population cost risk is
\begin{equation}
  R_C(h)=\E\left[C_{Y,h(X)}\right].
\end{equation}
For a test set $\{(x_i,y_i)\}_{i=1}^{n}$, the average test cost is
\begin{equation}
  \widehat{\operatorname{ATC}}(h)=
  \frac{1}{n}\sum_{i=1}^n C_{y_i,h(x_i)}.
\end{equation}
The normalized ATC used in the repository is
\begin{equation}
  \widehat{\operatorname{nATC}}(h)=
  \frac{\widehat{\operatorname{ATC}}(h)}
       {\max_{a\ne b} C_{a,b}},
\end{equation}
whenever the maximum off-diagonal cost is positive.
\end{definition}

\begin{definition}[Target-min recall]
Let $\mathcal{T}\subseteq[K]$ be the cared set of target classes.
The target-min recall is
\begin{equation}
  \operatorname{Recall}_{\min}(h;\mathcal{T})
  =
  \min_{y\in\mathcal{T}}
  P(h(X)=y\mid Y=y).
\end{equation}
PMI-10 uses the four cared NDC classes from the fixed profile.
\end{definition}

\begin{proposition}[Bayes rule for expected cost]
For any $x$, a cost-minimizing classifier satisfies
\begin{equation}
  h_C^\star(x)\in
  \argmin_{j\in[K]}
  \sum_{y=1}^{K}\eta_y(x)C_{y,j}.
\end{equation}
\end{proposition}

\begin{proof}
Condition on $X=x$.
If the classifier predicts class $j$, its conditional expected cost is
$\E[C_{Y,j}\mid X=x]=\sum_y P(Y=y\mid X=x)C_{y,j}$.
The global risk decomposes pointwise as
$R_C(h)=\E_X[\E[C_{Y,h(X)}\mid X]]$, so minimizing the inner expression for each
$x$ gives the stated rule.
\end{proof}

\begin{remark}[The argmax-training constraint]
The HPO rows discussed here use argmax inference rather than calibrated cost-min inference.
Thus, to reduce clean-test ATC, the cost matrix must change the learned score geometry itself:
\begin{equation}
  h_f(x)=\argmax_{k} f_k(x).
\end{equation}
This is why a loss-level explanation matters.
The relevant comparison is not ``which method has the Bayes cost-min decision rule,'' but which training procedure makes the argmax scores safer under the fixed PMI cost matrix.
\end{remark}

\section{NICME Formulation}

For one example with true label $y$ and raw logits $z=f_\theta(x)\in\R^K$, the repository implementation
\pathref{utils/loss\_functions.py} applies a positive pairwise cost offset to each non-target logit:
\begin{equation}
  z'_k
  =
  z_k+\alpha\,\I[k\ne y]\log\left(\max(C_{y,k},1)\right),
  \label{eq:nicme-adjusted-logits}
\end{equation}
where $\alpha=\code{nicme\_logit\_cost\_scale}$.
The NICME pairwise CE component is
\begin{align}
  L_{\mathrm{NICMECE}}(z,y)
  &=
  -\log
  \frac{\exp(z'_y)}{\sum_{j=1}^K \exp(z'_j)}
  \nonumber\\
  &=
  \log\left(
    e^{z_y}
    +
    \sum_{k\ne y} e^{z_k} C_{y,k}^{\alpha}
  \right)-z_y,
  \label{eq:nicme-ce}
\end{align}
where ordinary noncritical costs $C_{y,k}=1$ add no offset.

The hybrid loss adds a normalized expected-cost penalty on the original logits.
Let $p=\softmax(z)$ and $\tilde C=C/(C_{\max}+\epsilon)$, with
$C_{\max}=\max_{a,b}C_{a,b}$.
Then
\begin{equation}
  L_{\mathrm{NICME}}(z,y)
  =
  L_{\mathrm{NICMECE}}(z,y)
  +
  \lambda
  \sum_{k=1}^{K}\tilde C_{y,k}p_k,
  \label{eq:nicme-loss}
\end{equation}
up to the repository's optional warmup schedule for $\lambda=\code{cs\_lambda}$.

\subsection{Pairwise Cost Margins}

\begin{proposition}[Pairwise-margin theorem]
Fix an example with label $y$ and a wrong class $k\ne y$ with $C_{y,k}\ge 1$.
For any tolerance $\rho>0$, the odds contribution of class $k$ relative to the target in the NICME softmax satisfies
\begin{equation}
  \frac{\exp(z'_k)}{\exp(z'_y)}\le \rho
  \quad\Longleftrightarrow\quad
  z_y-z_k \ge \alpha\log C_{y,k}+\log(1/\rho).
  \label{eq:margin-equivalence}
\end{equation}
Thus a costly wrong class requires an extra raw-logit margin
$\alpha\log C_{y,k}$.
\end{proposition}

\begin{proof}
Since $z'_y=z_y$ and $z'_k=z_k+\alpha\log C_{y,k}$ for $k\ne y$,
\begin{equation}
  \frac{\exp(z'_k)}{\exp(z'_y)}
  =
  \exp(z_k-z_y)C_{y,k}^{\alpha}.
\end{equation}
The inequality $\exp(z_k-z_y)C_{y,k}^{\alpha}\le\rho$ is equivalent, after
taking logarithms, to
$z_y-z_k\ge \alpha\log C_{y,k}+\log(1/\rho)$.
\end{proof}

\begin{remark}[Why this remains active after correct classification]
The theorem does not condition on whether $\argmax_j z_j=y$.
The costly class is offset in the loss for every training example with label $y$.
Therefore NICME can keep increasing safety margins even after ordinary CE would already view the example as easy.
This directly addresses a terminal-phase-training concern raised by the CSADA paper: once overparameterized networks nearly interpolate the training set, losses that only react to observed mistakes can lose cost-sensitive signal \cite{chen2025csada}.
\end{remark}

\subsection{Expected-Cost Gradient}

\begin{proposition}[Expected-cost-gradient theorem]
Let
\begin{equation}
  R_y(z)=\sum_{k=1}^{K}\tilde C_{y,k}p_k,
  \qquad p=\softmax(z).
\end{equation}
Then for each logit coordinate $\ell$,
\begin{equation}
  \frac{\partial R_y(z)}{\partial z_\ell}
  =
  p_\ell
  \left(
    \tilde C_{y,\ell}
    -
    \sum_{k=1}^{K}\tilde C_{y,k}p_k
  \right).
  \label{eq:expected-cost-gradient}
\end{equation}
Consequently, the regularizer directly pushes down logits for classes whose cost is above the current expected cost under $p$.
\end{proposition}

\begin{proof}
The softmax derivative is
\begin{equation}
  \frac{\partial p_k}{\partial z_\ell}
  =
  p_k(\I[k=\ell]-p_\ell).
\end{equation}
Therefore
\begin{align}
  \frac{\partial R_y}{\partial z_\ell}
  &=
  \sum_k \tilde C_{y,k}
  p_k(\I[k=\ell]-p_\ell)
  \nonumber\\
  &=
  \tilde C_{y,\ell}p_\ell
  -
  p_\ell\sum_k \tilde C_{y,k}p_k
  \nonumber\\
  &=
  p_\ell
  \left(
    \tilde C_{y,\ell}
    -
    \sum_k \tilde C_{y,k}p_k
  \right).
\end{align}
Under gradient descent, a positive derivative decreases $z_\ell$.
Thus if class $\ell$ has above-current-mean cost for true label $y$, the expected-cost term applies direct downward pressure to that logit.
\end{proof}

\begin{lemma}[Direct descent on clean expected cost]
For the regularizer alone, the infinitesimal update
$z(t)=z-t\nabla R_y(z)$ satisfies
\begin{equation}
  \left.\frac{d}{dt}R_y(z(t))\right|_{t=0}
  =
  -\|\nabla R_y(z)\|_2^2\le 0.
\end{equation}
\end{lemma}

\begin{proof}
Apply the chain rule:
$\frac{d}{dt}R_y(z-t\nabla R_y(z))|_{t=0}
=\nabla R_y(z)^\top(-\nabla R_y(z))$.
\end{proof}

\keybox{
\textbf{Interpretation.}
NICME has two clean-data mechanisms.
The $\alpha$ term creates pairwise margins in raw logit space.
The $\lambda$ term directly reduces clean softmax mass on high-cost wrong classes.
Neither mechanism requires a generated adversarial example to become active.
}

\section{CSADA Formulation And Local-Limit Analysis}

CSADA was proposed specifically because ordinary cost-sensitive reweighting can be ineffective for overparameterized DNNs that interpolate the training set \cite{chen2025csada}.
The published method uses targeted adversarial examples to push decision boundaries in cost-aware directions on natural examples.
The repository implementation follows this spirit with a targeted PGD-style inner loop.

For each label $y$, define a selected critical target
\begin{equation}
  t(y)\in\argmax_{k} C_{y,k},
\end{equation}
and train on selected rows where $C_{y,t(y)}>1$.
Given an input $x$, the inner loop approximately solves a targeted attack:
\begin{equation}
  \delta^\star
  \approx
  \argmin_{\|\delta\|_\infty\le \epsilon}
  \operatorname{CE}(f_\theta(x+\delta),t(y)),
  \label{eq:csada-inner}
\end{equation}
which increases the model's score for the costly wrong class $t(y)$.
The outer loss is then
\begin{equation}
  L_{\mathrm{CSADA}}
  =
  \operatorname{CE}(f_\theta(x),y)
  +
  \lambda_{\mathrm{adv}}
  w_y
  \operatorname{CE}(f_\theta(x+\delta^\star),y),
  \label{eq:csada-outer}
\end{equation}
where $w_y$ is a normalized cost-derived weight.

\subsection{What CSADA Optimizes Locally}

Let
\begin{equation}
  m_t(x)=f_y(x)-f_t(x)
\end{equation}
be the clean raw margin between the true class and the targeted costly class.
Under a first-order expansion,
\begin{equation}
  f_t(x+\delta)-f_y(x+\delta)
  \approx
  f_t(x)-f_y(x)
  +
  \nabla_x(f_t-f_y)(x)^\top \delta.
\end{equation}
The best possible first-order increase of the wrong-minus-true margin inside an
$\ell_\infty$ ball is bounded by duality:
\begin{equation}
  \max_{\|\delta\|_\infty\le \epsilon}
  \nabla_x(f_t-f_y)(x)^\top\delta
  =
  \epsilon\|\nabla_x(f_t-f_y)(x)\|_1.
  \label{eq:dual-bound}
\end{equation}

\begin{proposition}[Local-CSADA limitation]
In the first-order regime, a targeted CSADA example can only expose a costly boundary for target
$t$ when the clean margin is small enough relative to local input sensitivity:
\begin{equation}
  m_t(x)
  \lesssim
  \epsilon\|\nabla_x(f_t-f_y)(x)\|_1.
  \label{eq:reachability}
\end{equation}
If the costly class has high clean softmax mass but the boundary is not reachable by the selected
perturbation path, or if another costly class is not selected as $t(y)$, the CSADA inner loop need not
produce a direct clean-logit penalty for that class.
\end{proposition}

\begin{proof}
The targeted inner loop tries to increase the relative score of $t$ against $y$.
Under the linear approximation, crossing the target-vs-true boundary requires
$f_t(x+\delta)-f_y(x+\delta)\ge 0$.
Using \eqref{eq:dual-bound}, the largest first-order increase is
$\epsilon\|\nabla_x(f_t-f_y)(x)\|_1$.
If this quantity is smaller than the clean margin $m_t(x)=f_y(x)-f_t(x)$, no perturbation in the local
linear $\ell_\infty$ model reaches that boundary.
The outer adversarial CE is then driven by a perturbed point whose effect on the clean logit vector is indirect.
The formulation also targets selected costly classes, not the entire row $C_{y,\cdot}$.
\end{proof}

\begin{remark}[Why this does not refute CSADA]
The proposition is a limitation, not a failure proof.
CSADA can be highly effective when the costly boundary is locally reachable and when boundary-supporting adversarial samples are the missing signal.
Indeed, Chen et al. report strong PMI improvements for CSADA over the baseline ResNet-34 pipeline \cite{chen2025csada}.
The point is narrower: CSADA's cost signal is mediated by local perturbation geometry, whereas NICME's cost signal is explicit in clean logit space.
\end{remark}

\subsection{Robust Optimization Tradeoff}

Adversarial training is usually formalized as robust optimization over an uncertainty set \cite{madry2018towards}.
The broader adversarial-robustness literature shows that robust learning can require more data \cite{schmidt2018robust}, can trade off with natural accuracy \cite{zhang2019trades}, and may suppress non-robust but predictive features \cite{ilyas2019features}.
Those facts do not imply that CSADA must reduce clean accuracy.
They do imply that, when the evaluation objective is clean-test ATC and balanced accuracy rather than adversarial robustness itself, adversarial augmentation can spend capacity on a surrogate robustness problem that is only indirectly aligned with the final metric.

\section{Why NICME Can Win On PMI-10}

The observed repository result is not mysterious under the above formulations.
PMI-10 is a clean-test, multiclass, pairwise-asymmetric cost problem.
The HPO-selected family uses argmax inference, so cost sensitivity must be encoded in the logits.
NICME is exactly a cost-conditioned logit-geometry method.

\begin{enumerate}
  \item \textbf{It separates pairwise costs from class imbalance.}
  Menon et al. show that logit adjustment can be statistically grounded as a margin mechanism for long-tail learning \cite{menon2021logit}.
  NICME changes the source of the margin from class priors to the pairwise cost matrix.
  That is better matched to PMI-10, where the important structure is not just ``rare class'' versus ``common class,'' but directed confusions such as true NDC $50111$-$0434$ predicted as $00591$-$0461$.
  \item \textbf{It keeps cost pressure active after correctness.}
  Ordinary CE becomes small once $z_y$ dominates.
  CSADA tries to recover signal by constructing boundary samples.
  NICME instead makes the clean loss demand a larger margin against costly wrong classes by construction.
  \item \textbf{It directly optimizes a differentiable proxy for clean ATC.}
  The expected-cost regularizer is a smooth surrogate for $\sum_k C_{y,k}\I[h(x)=k]$.
  Modern cost-sensitive work continues to emphasize that expected misclassification cost, rather than error rate alone, is the measure of interest in many applications \cite{komisarenko2025cost}.
  \item \textbf{It avoids some adversarial-training clean-accuracy costs.}
  CSADA can move boundaries in useful cost-aware directions, but its perturbation mechanism can also change the learned feature preference.
  In this run, the empirical fingerprint is consistent with that risk: CSADA attains excellent cared-class recall but much lower balanced accuracy than NICME.
\end{enumerate}

\section{Empirical Comparison}

The current repository comparison defines ``SOTA'' as the completed internal PMI-10 baseline suite, not a global external benchmark.
The strongest previous recall-first baseline row was CSADA in the LR $5\mathrm{e}{-5}$ profile.
The best observed NICME HPO test row is run 95 with $\alpha=0.5$ and $\lambda=0.07$.

\begin{center}
\small
\begin{tabular}{>{\raggedright\arraybackslash}p{0.32\linewidth}rrr}
\toprule
Metric & CSADA baseline & NICME run 95 & Direction \\
\midrule
Target-min recall & 0.96875 & 0.96875 & tie \\
Target-macro recall & 0.984375 & 0.984375 & tie \\
Normalized ATC & 0.015625 & 0.0046875 & NICME lower by 70.0\% \\
ATC & 0.15625 & 0.046875 & NICME lower by 70.0\% \\
Balanced accuracy & 0.871875 & 0.98125 & NICME higher \\
Macro F1 & 0.851092 & 0.981386 & NICME higher \\
Critical-pair errors & 1 & 1 & tie \\
\bottomrule
\end{tabular}
\end{center}

\noindent Repository sources:
\begin{itemize}
  \item \pathref{results/pmi10 historical alpha-lambda comparison summary}
  \item \pathref{results/pmi10 historical alpha-lambda grid ledger}
  \item \pathref{results/pmi10\_sota\_pretty\_balanced\_triple\_lr\_20260503/comparison\_summary.md}
\end{itemize}

\keybox{
\textbf{Best supported empirical claim.}
NICME run 95 is the best observed recall-first cost-sensitive tradeoff among the completed repository baselines and completed HPO rows:
it matches CSADA's cared-class recall while substantially reducing normalized ATC and improving balanced accuracy.
}

\section{Caveats And Falsification Tests}

\subsection{Caveats}

\begin{itemize}
  \item \textbf{Run 95 is test-best, not validation-selected.}
  The validation-selected HPO row is run 89.
  Run 95 is validation rank 2 and is the best observed test row.
  A final paper claim should use a prospective selection rule or fresh confirmation.
  \item \textbf{Run 69 ties the main cost-sensitive metrics.}
  Run 95 beats run 69 only by a very small macro-F1 margin.
  The evidence supports a stable high-performing region near moderate $\alpha$ and $\lambda$, not a magical single point.
  \item \textbf{NICME does not minimize every metric.}
  If the sole objective is minimum normalized ATC or zero critical errors, \code{ce\_cost\_min\_inference\_pretty} at LR $1\mathrm{e}{-4}$ remains better in the current table, with normalized ATC $0.0025$ and zero critical errors, but target-min recall $0.90625$.
  \item \textbf{CSADA has a legitimate theoretical niche.}
  CSADA was designed for overparameterized models whose training examples are already fit too well for ordinary cost weighting to matter.
  When local boundary augmentation is the missing signal, CSADA may outperform a clean-logit method.
\end{itemize}

\subsection{Falsification Tests}

The theory above would be weakened if any of the following occurred in a prospective experiment:

\begin{enumerate}
  \item A multi-seed nested-validation run selects CSADA or AP-CSADA with equal recall and lower clean normalized ATC.
  \item NICME's advantage disappears when run 95 is not allowed to be selected after inspecting test metrics.
  \item Ablating the expected-cost regularizer leaves performance unchanged, contradicting the probability-mass explanation.
  \item Ablating the pairwise logit margin leaves performance unchanged, contradicting the margin explanation.
  \item Measuring margins shows no larger $z_y-z_k$ gaps on high-cost PMI pairs for NICME than for CSADA.
\end{enumerate}

\section{Conclusion}

NICME and CSADA solve related but not identical optimization problems.
CSADA uses adversarial data augmentation to create boundary-supporting samples for critical pairs, a sensible response to overparameterized interpolation.
NICME instead changes the clean training objective so that high-cost wrong classes require larger raw margins and receive direct probability-mass penalties.
For a clean-test, argmax-inference, recall-first PMI-10 evaluation, that alignment gives a principled reason for the observed result:
NICME can preserve CSADA's cared-class recall while reducing expected cost and avoiding the clean balanced-accuracy degradation seen in the CSADA row.

\begin{thebibliography}{11}

\bibitem{bartlett2006convexity}
Peter L. Bartlett, Michael I. Jordan, and Jon D. McAuliffe.
\newblock Convexity, classification, and risk bounds.
\newblock \emph{Journal of the American Statistical Association}, 101(473):138--156, 2006.
\newblock \url{https://doi.org/10.1198/016214505000000907}

\bibitem{chen2025csada}
Qiyuan Chen, Raed Al Kontar, Maher Nouiehed, X. Jessie Yang, and Corey Lester.
\newblock Rethinking cost-sensitive classification in deep learning via adversarial data augmentation.
\newblock \emph{INFORMS Journal on Data Science}, 4(1):1--19, 2025.
\newblock \url{https://doi.org/10.1287/ijds.2022.0033}

\bibitem{elkan2001foundations}
Charles Elkan.
\newblock The foundations of cost-sensitive learning.
\newblock In \emph{Proceedings of the Seventeenth International Joint Conference on Artificial Intelligence}, 2001.
\newblock \url{https://cseweb.ucsd.edu/~elkan/rescale.pdf}

\bibitem{guo2017calibration}
Chuan Guo, Geoff Pleiss, Yu Sun, and Kilian Q. Weinberger.
\newblock On calibration of modern neural networks.
\newblock In \emph{Proceedings of the 34th International Conference on Machine Learning}, 2017.
\newblock \url{https://proceedings.mlr.press/v70/guo17a.html}

\bibitem{ilyas2019features}
Andrew Ilyas, Shibani Santurkar, Dimitris Tsipras, Logan Engstrom, Brandon Tran, and Aleksander Madry.
\newblock Adversarial examples are not bugs, they are features.
\newblock In \emph{Advances in Neural Information Processing Systems}, 2019.
\newblock \url{https://papers.nips.cc/paper/8307-adversarial-examples-are-not-bugs-they-are-features}

\bibitem{komisarenko2025cost}
Viacheslav Komisarenko and Meelis Kull.
\newblock Cost-sensitive classification with cost uncertainty: do we need surrogate losses?
\newblock \emph{Machine Learning}, 114:132, 2025.
\newblock \url{https://doi.org/10.1007/s10994-024-06634-8}

\bibitem{madry2018towards}
Aleksander Madry, Aleksandar Makelov, Ludwig Schmidt, Dimitris Tsipras, and Adrian Vladu.
\newblock Towards deep learning models resistant to adversarial attacks.
\newblock \emph{International Conference on Learning Representations}, 2018.
\newblock \url{https://arxiv.org/abs/1706.06083}

\bibitem{menon2021logit}
Aditya Krishna Menon, Sadeep Jayasumana, Ankit Singh Rawat, Himanshu Jain, Andreas Veit, and Sanjiv Kumar.
\newblock Long-tail learning via logit adjustment.
\newblock \emph{International Conference on Learning Representations}, 2021.
\newblock \url{https://openreview.net/forum?id=37nvvqkCo5}

\bibitem{schmidt2018robust}
Ludwig Schmidt, Shibani Santurkar, Dimitris Tsipras, Kunal Talwar, and Aleksander Madry.
\newblock Adversarially robust generalization requires more data.
\newblock In \emph{Advances in Neural Information Processing Systems}, 2018.
\newblock \url{https://arxiv.org/abs/1804.11285}

\bibitem{xin2025provably}
Yuan Xin, Dingfan Chen, Michael Backes, and Xiao Zhang.
\newblock Provably cost-sensitive adversarial defense via randomized smoothing.
\newblock In \emph{Proceedings of the 42nd International Conference on Machine Learning}, 2025.
\newblock \url{https://proceedings.mlr.press/v267/xin25a.html}

\bibitem{zhang2019trades}
Hongyang Zhang, Yaodong Yu, Jiantao Jiao, Eric Xing, Laurent El Ghaoui, and Michael Jordan.
\newblock Theoretically principled trade-off between robustness and accuracy.
\newblock In \emph{Proceedings of the 36th International Conference on Machine Learning}, 2019.
\newblock \url{https://arxiv.org/abs/1901.08573}

\end{thebibliography}

\end{document}
